\documentclass[a4paper,12pt]{article}
\usepackage[utf8]{inputenc}
\usepackage[T1]{fontenc}
\usepackage[ngerman]{babel}
\usepackage{amsmath}
\usepackage{amssymb}
\usepackage{enumitem}
\usepackage{geometry}
\geometry{a4paper, margin=2.5cm}
\usepackage{parskip}

\title{Warum in Aufgabe 3\,b) die Gleichheit nur bei gleichen Koordinaten gilt\\[0.4em]
\large 61211 Analysis -- Probeklausur}
\author{}
\date{}

\begin{document}
\maketitle

\section*{Frage}
Muss bei Gleichheit in Aufgabe~3\,a) wirklich in allen Koordinaten derselbe Wert stehen,
also $x_1=x_2=\dots=x_n$? Kann man das nicht mit der euklidischen Norm untersuchen und
dann auch ungleiche Koordinaten zulassen?

\textbf{Antwort:} Ja, es muss $x_1=\dots=x_n$ gelten. Der Zugang über die euklidische Norm
ist völlig richtig -- er liefert aber genau dasselbe Ergebnis.

\section*{Ausgangslage}
Betrachtet werden alle Vektoren $x=(x_1,\dots,x_n)\in\mathbb{R}^n$ mit
\[
  x_1+x_2+\dots+x_n=c .
\]
Die Behauptung aus a) lautet
\[
  \sum_{k=1}^n x_k^2 \;\ge\; \frac{c^2}{n}.
\]

\section*{Weg 1: Gleichheitsfall der Cauchy--Schwarzschen Ungleichung}
In der Cauchy--Schwarzschen Ungleichung
\[
  \Big|\sum_{k=1}^n x_k\,y_k\Big|\le \Big(\sum_{k=1}^n x_k^2\Big)^{1/2}\Big(\sum_{k=1}^n y_k^2\Big)^{1/2}
\]
gilt Gleichheit \emph{genau dann}, wenn $x$ und $y$ linear abhängig sind. In a) ist
$y=(1,1,\dots,1)$. Linear abhängig von $(1,\dots,1)$ bedeutet
\[
  x_k=\lambda \qquad\text{für alle } k,
\]
also gleiche Koordinaten. Einsetzen in die Nebenbedingung liefert $n\lambda=c$ und damit
\[
  \lambda=\frac{c}{n}
  \qquad\Longrightarrow\qquad
  x=\Big(\tfrac{c}{n},\tfrac{c}{n},\dots,\tfrac{c}{n}\Big).
\]
Es gibt also genau \emph{einen} Vektor mit Gleichheit.

\section*{Weg 2: Die euklidische (geometrische) Sicht}
$\sum_{k=1}^n x_k^2=\lVert x\rVert_2^2$ zu minimieren heißt, auf der Hyperebene
\[
  H=\{\,x\in\mathbb{R}^n : x_1+\dots+x_n=c\,\}
\]
den Punkt mit kleinstem Abstand zum Ursprung zu suchen (den Lotfußpunkt). Die
Normalenrichtung von $H$ ist $(1,\dots,1)$; der Lotfußpunkt liegt daher auf dieser Richtung:
\[
  x=\frac{c}{n}\,(1,\dots,1).
\]
Da $\lVert x\rVert_2^2$ \emph{strikt konvex} ist, ist dieser nächstgelegene Punkt
\emph{eindeutig} -- kein zweiter Vektor erreicht denselben Minimalwert. Genau das ist die
Gleichheit in a).

\section*{Konkreter Test ($n=2$, $c=2$)}
Alle zulässigen $x$ haben die Form $x=(1+t,\,1-t)$ mit $t\in\mathbb{R}$, denn dann ist
$x_1+x_2=2$. Es folgt
\[
  x_1^2+x_2^2=(1+t)^2+(1-t)^2=2+2t^2\;\ge\;2=\frac{c^2}{n}.
\]
\begin{itemize}[itemsep=2pt]
  \item $t=0\ \Rightarrow\ x=(1,1)\ \Rightarrow\ $ Wert $2$ \;(das ist die Schranke $c^2/n$),
  \item $t=\tfrac12\ \Rightarrow\ x=(1{,}5,\,0{,}5)\ \Rightarrow\ $ Wert $2{,}5$ \;(größer),
  \item $t=1\ \Rightarrow\ x=(2,\,0)\ \Rightarrow\ $ Wert $4$ \;(noch größer).
\end{itemize}
Ungleiche Koordinaten sind also erlaubt, liefern aber \emph{stets} einen echt größeren Wert.
Der kleinstmögliche Wert (die Gleichheit) tritt nur bei gleichen Koordinaten auf.

\section*{Kernpunkt (auch per Lagrange)}
Auf der Hyperebene darf $x$ frei gewählt werden. Sobald man aber das Minimum von
$\sum x_k^2$ -- also die Gleichheit in a) -- verlangt, ist die Wahl erzwungen. Mit einem
Lagrange-Multiplikator $\mu$ für die Nebenbedingung $g(x)=\sum x_k-c=0$:
\[
  \nabla\Big(\sum_{k=1}^n x_k^2\Big)=\mu\,\nabla g(x)
  \quad\Longleftrightarrow\quad
  2x_k=\mu\ \text{ für alle } k
  \quad\Longrightarrow\quad
  x_1=x_2=\dots=x_n.
\]
Zusammen mit $\sum x_k=c$ ergibt sich erneut $x_k=\tfrac{c}{n}$.

\medskip
\noindent\textbf{Fazit:} Die Musterlösung zu 3\,b) ist korrekt: Gleichheit gilt genau für
$x=\big(\tfrac{c}{n},\dots,\tfrac{c}{n}\big)$.

\end{document}
